\subsection{The decomposition theorem}

THEOREM 1 (decomposition theorem). --- Let \((x_i)_{1 \leq i \leq p}\) and \((y_j)_{1 \leq j \leq q}\) be two finite sequences of positive elements of a lattice ordered group \(G\) such that \(\sum_{i=1}^{p} x_i = \sum_{j=1}^{q} y_j\) ; then there exists a double sequence \((z_{ij})_{1 \leq i \leq p, 1 \leq j \leq q}\) of positive elements of \(G\) such that \(x_i = \sum_{j=1}^{q} z_{ij}\) for all \(i\) , and \(y_j = \sum_{i=1}^{p} z_{ij}\) for all \(j\) .

\begin{enumerate}
\def\labelenumi{\arabic{enumi})}
\tightlist
\item
  Let us prove the theorem first for the case \(p = q = 2\) . Let \(x, x', y, y'\) be positive elements of \(G\) such that \(x + x' = y + y'\) , and put \(a = \sup (0, x - y')\) . Since
\end{enumerate}

\[
x - y ^ {\prime} = y - x ^ {\prime}
\]

is less than \(x\) and less than \(y\) , it follows that \(b = x - a\) and \(c = y - a\) are positive, as is \(d = a - (x - y')\) . Also we have

\[
x = a + b, x ^ {\prime} = c + d, y = a + c \quad \text { and } \quad y ^ {\prime} = b + d.
\]

\begin{enumerate}
\def\labelenumi{\arabic{enumi})}
\setcounter{enumi}{1}
\tightlist
\item
  Let us now show that if the theorem is true for \(p < m\) and \(q = n\) ( \(m > 2\) , \(n \geqslant 2\) ) then it is true for \(p = m\) and \(q = n\) . By hypothesis we have \(x_{m} + \sum_{i=1}^{m-1} x_{i} = \sum_{j=1}^{n} y_{j}\) . The theorem being true for \(p = 2\) and \(q = n\) , there exist two finite sequences \((z_{j}')\) , \((z_{j}''')\) of \(n\) positive terms such that \(\sum_{i=1}^{m-1} x_{i} = \sum_{j=1}^{n} z_{j}'\) , \(x_{m} = \sum_{j=1}^{n} z_{j}''\) , and \(y_{j} = z_{j}' + z_{j}''\) for \(1 \leqslant j \leqslant n\) . On the other hand, since the theorem is true for \(p = m-1\) and \(q = n\) , there exists a double sequence \((u_{ij})_{1 \leqslant i \leqslant m-1, 1 \leqslant j \leqslant n}\) such that \(x_{i} = \sum_{j=1}^{n} u_{ij}\) for \(1 \leqslant i \leqslant m-1\) , and \(z_{j}' = \sum_{i=1}^{m-1} u_{ij}\) for \(1 \leqslant j \leqslant n\) . Putting
\end{enumerate}

\[
z _ {i j} = u _ {i j} \quad \text { for } \quad 1 \leqslant i \leqslant m - 1, \quad \text { and } \quad z _ {m j} = z _ {j} ^ {\prime \prime} \quad (1 \leqslant j \leqslant n),
\]

we certainly obtain a double sequence satisfying the conditions of the theorem. 3) By interchanging the \(x_{i}\) and the \(y_{j}\) we see in the same way that, if the theorem is true for \(p = m\) and \(q < n\) ( \(m \geqslant 2\) , \(n > 2\) ), then it is true for \(p = m\) and \(q = n\) . The theorem is thus proved by double induction starting from the case \(p = q = 2\) , for it is trivially true whenever \(p \leqslant 1\) or \(q \leqslant 1\) .

COROLLARY. --- Let \(y, x_1, x_2, \ldots, x_n\) be \(n + 1\) positive elements of \(G\) such that \(y \leqslant \sum_{i=1}^{n} x_i\) ; then there exist \(n\) positive elements \(y_i (1 \leqslant i \leqslant n)\) such that \(y_i \leqslant x_i\) and \(y = \sum_{i=1}^{n} y_i\) .

It is sufficient to apply theorem 1 to the sequence \((x_{i})\) and the sequence consisting of the two elements \(y\) and \(z = \left(\sum_{i=1}^{n} x_{i}\right) - y\) .

